\usepackage{enumitem}
\setenumerate[0]{label=\arabic*)} % Изменение вида нумерации списков
\renewcommand{\labelitemi}{---}

%==========Для подключения русского языка=====================
\usepackage{fontspec} % Подключаем пакет для работы со шрифтами
\setmainfont{Times New Roman} % указывает XeLaTeX какой шрифт использовать
\usepackage[utf8]{inputenc}     % Поддержка UTF-8
%=============================================================


\usepackage{caption}
\captionsetup[table]{justification=raggedright,singlelinecheck=off} % Изменение подписей к таблицам
\captionsetup{labelsep=endash, justification=centering} % Настройка подписей float объектов
\captionsetup[figure]{name=Рисунок} % Изменение подписей к рисункам

\usepackage{geometry}
\geometry{a4paper,left=3.1cm,right=1.1cm,top=2cm,bottom=2cm,bindingoffset=0cm}
%\geometry{left=30mm}
%\geometry{right=15mm}
%\geometry{top=20mm}
%\geometry{bottom=20mm}
\setlength{\parindent}{1.25cm}
%\linespread{1.5}

% --- Пакеты для красивого оформления алгоритмов ---
\usepackage{algorithm}
\usepackage{algorithmicx}
\usepackage{algpseudocode}

%===========ОПЦИОНАЛЬНО==========

\addbibresource{biblio.bib}

%\usepackage[T1]{fontenc}
%\usepackage{lmodern}

\usepackage{multirow}
\usepackage{makecell}

\usepackage{ulem} % Нормальное нижнее подчеркивание
% Дополнительное окружения для подписей
\usepackage{array}
\newenvironment{signstabular}[1][1]{
	\renewcommand*{\arraystretch}{#1}
	\tabular
}{
	\endtabular
}
\def\arraybackslash{\let\\\tabularnewline}

\usepackage{pdfpages}
\usepackage{pgfplots}
\pgfplotsset{compat=1.9}
\usetikzlibrary{datavisualization}
\usetikzlibrary{datavisualization.formats.functions}
\usepackage{csvsimple} % генерим таблички из csv

\usepackage{graphicx}
\newcommand{\imgHeight}[3] {
	\begin{figure}[H]
		\center{\includegraphics[height=#1]{inc/img/#2}}
		\caption{#3}
		\label{img:#2}
	\end{figure}
}

\newcommand{\imgWid}[3] {
	\begin{figure}[H]
		\center{\includegraphics[width=#1]{inc/img/#2}}
		\caption{#3}
		\label{img:#2}
	\end{figure}
}

\newcommand{\imgScale}[3] {
	\begin{figure}[h!]
		\center{\includegraphics[scale=#1]{inc/img/#2}}
		\caption{#3}
		\label{img:#2}
	\end{figure}
}

\usepackage{listingsutf8}
\lstset{columns=fixed,basicstyle=\small,breaklines=true,inputencoding=utf8,keywordstyle=\bfseries\underbar,frame=single,tabsize=2,xleftmargin=20pt,xrightmargin=5pt}
%\lstset{language=sql}
\lstset{numbers=left,numberstyle=\small}

\lstset{extendedchars=\true}

% \usepackage{listings}
% \usepackage{xcolor}
% \lstset{ %
% 	basicstyle=\small\sffamily,			% размер и начертание шрифта для подсветки кода
% 	numbers=left,						% где поставить нумерацию строк (слева\справа)
% 	stepnumber=1,						% размер шага между двумя номерами строк
% 	numbersep=5pt,						% как далеко отстоят номера строк от подсвечиваемого кода
% 	frame=single,						% рисовать рамку вокруг кода
% 	tabsize=4,							% размер табуляции по умолчанию равен 4 пробелам
% 	captionpos=t,						% позиция заголовка вверху [t] или внизу [b]
% 	breaklines=true,					
% 	breakatwhitespace=true,				% переносить строки только если есть пробел
% 	escapeinside={\#*}{*)},				% если нужно добавить комментарии в коде
% 	backgroundcolor=\color{white},
% }
% Команда для римских цифр
%\newcommand{\rom}[1]{\MakeUppercase{\romannumeral #1}}
\usepackage{placeins}

\usepackage{xcolor}
\usepackage{listings}

% --- Настройка стиля листингов ---
% --- Настройка стиля листингов ---
\lstset{
	language=C++,
%	basicstyle=\footnotesize\fontfamily{ccr}\selectfont, % ccr = Courier
	basicstyle=\footnotesize\ttfamily,
	numbers=left,
	numberstyle=\tiny\color{gray},
	frame=tb,
	tabsize=4,
	showstringspaces=false,
	breaklines=true,
	keywordstyle=\color{blue},
	commentstyle=\color{green!40!black},
	stringstyle=\color{purple},
	backgroundcolor=\color{black!5},
	extendedchars=true,  % Включаем поддержку расширенных символов
	literate= {А}{{\CYRA}}1 {Б}{{\CYRB}}1 {В}{{\CYRV}}1 {Г}{{\CYRG}}1 {Д}{{\CYRD}}1 {Е}{{\CYRE}}1 {Ё}{{\CYRYO}}1
	{Ж}{{\CYRZH}}1 {З}{{\CYRZ}}1 {И}{{\CYRI}}1 {Й}{{\CYRI breve}}1 {К}{{\CYRK}}1 {Л}{{\CYRL}}1 {М}{{\CYRM}}1
	{Н}{{\CYRN}}1 {О}{{\CYRO}}1 {П}{{\CYRP}}1 {Р}{{\CYRR}}1 {С}{{\CYRS}}1 {Т}{{\CYRT}}1 {У}{{\CYRU}}1
	{Ф}{{\CYRF}}1 {Х}{{\CYRH}}1 {Ц}{{\CYRC}}1 {Ч}{{\CYRCH}}1 {Ш}{{\CYRSH}}1 {Щ}{{\CYRSHCH}}1
	{Ъ}{{\CYRHardSign}}1 {Ы}{{\CYRY}}1 {Ь}{{\CYRSoftSign}}1 {Э}{{\CYRE}}1 {Ю}{{\CYRYU}}1 {Я}{{\CYRYA}}1
	{а}{{\cyra}}1 {б}{{\cyrb}}1 {в}{{\cyrv}}1 {г}{{\cyrg}}1 {д}{{\cyrd}}1 {е}{{\cyre}}1 {ё}{{\cyryo}}1
	{ж}{{\cyrzh}}1 {з}{{\cyrz}}1 {и}{{\cyri}}1 {й}{{\cyri breve}}1 {к}{{\cyrk}}1 {л}{{\cyrl}}1 {м}{{\cyrm}}1
	{н}{{\cyrn}}1 {о}{{\cyro}}1 {п}{{\cyrp}}1 {р}{{\cyrr}}1 {с}{{\cyrs}}1 {т}{{\cyrt}}1 {у}{{\cyru}}1
	{ф}{{\cyrf}}1 {х}{{\cyrh}}1 {ц}{{\cyrc}}1 {ч}{{\cyrch}}1 {ш}{{\cyrsh}}1 {щ}{{\cyrshch}}1
	{ъ}{{\cyrhard}}1 {ы}{{\cyry}}1 {ь}{{\cyrsoft}}1 {э}{{\cyre}}1 {ю}{{\cyryu}}1 {я}{{\cyrya}}1
}